\scanpage{9}
\mspage{7}
Ces structures sont héritées par les
\[
\boxed{\begin{aligned}
C^{\bullet\bullet}_{\mathrm{DRpd}}(\unsure{\mathscr X|X|\mathcal X}_{\ast},k)&\overset{\mathrm{déf}}{=}C^\bullet_{\mathrm{DRpd}}(\mathscr X_{\ast};k\{T\},k\{T\}^{+},T)\\
&=\mathrm{Hom}\big(\mathscr X_\ast,C^{\bullet\bullet}_{\mathrm{DRpd}\,\ast}(k)\big)
\end{aligned}}
\]
et dépendent de façon contravariante de $\mathscr X_{\ast}$, (covariante de $k$).

$\Phi_{k\ast}$ est \deleted{\unsure{$k$}}\inserted{un $k$-Module $1/2$ simplicial} homotope à 0 donc $\unsure{\Gamma|\mathscr T}^{p}_k(\Phi_{k\ast})\otimes\bigwedge^q\Psi_{k\ast}$ \deleted{\unsure{est}} est homotope à $\Gamma^p_k(0)\otimes\bigwedge^q\Psi_{k\ast}$, donc hom. à 0 si $p\ne 0$.

Donc $C^{p,q}_{\mathrm{DRpd}\,\ast}(k)$ est hom à 0 (donc \deleted{flasque}\inserted{acyclique}) si $p=0$. En degré total donné \inserted{$n$}, on a
\diagram{\begin{tikzcd}[ampersand replacement=\&,column sep=tiny,row sep=small]
0 \arrow[r] \& C^{n,0}_{\ast} \arrow[r] \& C^{n-1,1}_{\ast} \arrow[r] \& C^{n-2,2}_{\ast} \arrow[r,dash] \& \cdots \& C^{\unsure{1|0},n-1}_{\ast} \arrow[r,phantom,"\unsure{\text{---}|\leftarrow|\to}"] \& C^{0,n}_{\ast} \arrow[r] \& 0 \arrow[d] \\
\& k_\ast \arrow[u] \& \& \& \& \& \& 0 \arrow[d] \\
\& \& \& \& \& \& \& 0
\end{tikzcd}}
i.e. on trouve une résolution \inserted{de long. $n$} \deleted{flasque}\inserted{de} $k_{\ast}$ par des \deleted{groupes}\inserted{$k$-Modules} $1/2$ simpliciaux qui sont acycliques sauf le dernier -- donc on peut la considérer comme \unsure{une|un} tronqué de \deleted{l'indice}\inserted{degré} $n$ d'une résolution flasque de $k_{\ast}$ -- la cohomologie de ses sections sur un $\mathscr X_{\ast}$ est donc la coh. de $\mathscr X_{\ast}$ tronquée en degré $n$ :
\[
H_{\mathrm{DRpd}}^{p,q}(\mathscr X_{\ast},k)\simeq
\begin{cases}H^q(\mathscr X_{\ast},k)&\text{si }q\leq p+q\ \text{i.e. }p\geq0\\
0&\text{si }p<0
\end{cases}
\]

% archive sheet number: 8
